% 대재해 예방투자의 손실통제 효과에 관한 연구 — 발표 슬라이드
% 컴파일: xelatex catastrophe-prevention-investment-slides.tex (2회)
\documentclass[aspectratio=169,10pt]{beamer}

\usetheme[progressbar=frametitle,numbering=fraction,block=fill]{metropolis}
\usepackage{kotex}
\setmainhangulfont{NanumBarunGothic}[BoldFont={NanumBarunGothic Bold}]
\setsanshangulfont{NanumBarunGothic}[BoldFont={NanumBarunGothic Bold}]
\setmainfont{NanumBarunGothic}[BoldFont={NanumBarunGothic Bold}]
\setsansfont{NanumBarunGothic}[BoldFont={NanumBarunGothic Bold}]
\setsanshanjafont{WenQuanYi Zen Hei}
\setmainhanjafont{WenQuanYi Zen Hei}
\newfontfamily\fbfont{DejaVu Sans}
\newcommand{\fb}[1]{{\fbfont #1}}
\metroset{sectionpage=none}

\usepackage{booktabs}
\usepackage{colortbl}
\usepackage{tabularx}
\usepackage{array}
\usepackage{amsmath,amssymb}
\usepackage{pgfplots}
\pgfplotsset{compat=1.18}
\usetikzlibrary{arrows.meta,positioning,shapes.geometric}

\definecolor{domestic}{HTML}{2A78D6}
\definecolor{foreign}{HTML}{EB6834}
\definecolor{accent}{HTML}{C0392B}
\definecolor{muted}{HTML}{6B6B6B}
\definecolor{rowhl}{HTML}{FDEBD0}
\setbeamercolor{alerted text}{fg=accent}

\newcommand{\src}[1]{\par\vfill{\tiny\color{muted}#1}}
\newcolumntype{Y}{>{\raggedright\arraybackslash}X}
\renewcommand{\arraystretch}{1.15}

\title{대재해 예방투자의 손실통제 효과에 관한 연구}
\subtitle{A Study on the Loss-Control Effects of Catastrophe Prevention Investment}
\author{홍지민 \and 이항석 \and 정힘찬}
\date{}

\begin{document}

% ------------------------------------------------------------------
\maketitle

\begin{frame}{발표 순서}
  \tableofcontents
\end{frame}

% ==================================================================
\section{서론}

\begin{frame}{연구 배경}
  \begin{itemize}
    \item 기후변화 $\cdot$ 도시 집적 $\cdot$ 시설물 노후화 $\Rightarrow$ 대재해의 \alert{빈도와 피해 규모} 동시 증가
    \begin{itemize}
      \item 2023 튀르키예-시리아 지진: 사망자 약 6만 명
      \item 2025 영남권 초대형 산불: 정부 수립 이후 최대 산불 피해(복구비 1조 8,809억 원)
    \end{itemize}
    \item 대재해는 사후 대응$\cdot$복구만으로 감당하기 어려운 인적$\cdot$물적 손실을 야기
    \item[$\Rightarrow$] 사전 \alert{예방투자의 효과를 계량적으로 검증}하는 것이 재난관리 정책의 핵심 과제
    \item 그러나 예방투자 효과를 \textbf{빈도}와 \textbf{심도}로 분해하여 이론화하고\\
          국내 자료로 실증한 연구는 드묾
  \end{itemize}
\end{frame}

\begin{frame}{연구 목적과 구성}
  \begin{columns}[T]
    \begin{column}{0.32\textwidth}
      \begin{block}{① 사례조사 (Ⅱ장)}
        2021--2025년 국내외 주요 대재해 \textbf{14건}의 원인구조 비교\\[3pt]
        \small 예방투자의 공백이 재난 대형화에 미친 영향
      \end{block}
    \end{column}
    \begin{column}{0.32\textwidth}
      \begin{block}{② 이론모형 (Ⅲ장)}
        복합 포아송 손실모형\\[3pt]
        \small 빈도저감형 vs.\ 심도저감형 투자의 사회적 가치 비교
      \end{block}
    \end{column}
    \begin{column}{0.32\textwidth}
      \begin{block}{③ 실증분석 (Ⅳ장)}
        광역지자체 17개 $\times$ 2015--2024 패널\\[3pt]
        \small 재난방재 예산과 자연재해 피해의 빈도$\cdot$심도 관계
      \end{block}
    \end{column}
  \end{columns}
  \vspace{1em}
  \begin{center}
    핵심 질문: 예방투자는 재해의 \textbf{발생 자체}를 줄이는가,\\
    아니면 \alert{대형피해로의 발전}을 억제하는가?
  \end{center}
\end{frame}

\begin{frame}{선행연구}
  \begin{itemize}
    \item \textbf{손실통제 이론}
    \begin{itemize}\small
      \item Ehrlich \& Becker (1972): 자가보험(self-insurance)과 자가보호(self-protection) 구분
      \item Lee (1998): 손실확률$\cdot$손실규모에 동시 작용하는 SICP(self-insurance-cum-protection)
      \item Wong (2016), 홍지민 (2020, 2022): 2기간 모형으로 확장
      \item 홍순구 (1999), 이원돈$\cdot$Ewedemi (2009): 순수$\cdot$투자위험 공존, 자가보험 선호
    \end{itemize}
    \item \textbf{자연재해 실증}
    \begin{itemize}\small
      \item 이동근 외 (2017): 기후변화 시나리오별 산사태 위험지역$\cdot$피해규모 추정
      \item 배은경 외 (2018): 극단치이론을 이용한 한$\cdot$일$\cdot$중 지진위험 모형화
    \end{itemize}
    \item \textbf{본 연구의 차별점}: 예방투자 효과를 \alert{빈도저감}과 \alert{심도저감}의 두 경로로 구분하여 비교
  \end{itemize}
\end{frame}

% ==================================================================
\section{국내외 주요 대재해 사례조사}

\begin{frame}{분석 대상 사례 총괄 (1): 국내 7건}
  \footnotesize
  \begin{tabularx}{\textwidth}{@{}l l l Y r@{}}
    \toprule
    사례명 & 시점 & 유형 & 직접 원인 & 사망자 \\
    \midrule
    광주 화정 아파트 붕괴 & 2022.1. & 붕괴 & 무단 공법 변경, 동바리 조기 해체, 콘크리트 강도 미달 & 6명 \\
    이태원 다중밀집 참사 & 2022.10. & 사회재난 & 군중 밀집도 관리 및 예방적 인파 통제 부재 & 159명 \\
    오송 지하차도 침수 & 2023.7. & 풍수해 & 미호강 임시제방 부실, 교통통제 미실시 & 14명 \\
    화성 아리셀 공장 화재 & 2024.6. & 화재(산업) & 리튬 일차전지 열폭주, 안전교육$\cdot$대피체계 부재 & 23명 \\
    부천 호텔 화재 & 2024.8. & 화재(건축) & 객실 에어컨 발화, 스프링클러 미설치 & 7명 \\
    무안 제주항공 참사 & 2024.12. & 사회재난 & 조류충돌 후 동체착륙, 콘크리트 방위각시설 충돌 & 179명 \\
    영남권 초대형 산불 & 2025.3. & 산불 & 성묘객 실화 등 + 고온건조$\cdot$초속 25m 강풍 & 32명 \\
    \bottomrule
  \end{tabularx}
  \src{자료: 감사원(2023), 고용노동부(2024), 국토교통부 건설사고조사위원회(2022), 국회 국정조사특별위원회(2023), 중앙재난안전대책본부(2025), 항공$\cdot$철도사고조사위원회(2025)}
\end{frame}

\begin{frame}{분석 대상 사례 총괄 (2): 해외 7건}
  \footnotesize
  \begin{tabularx}{\textwidth}{@{}l l l Y r@{}}
    \toprule
    사례명 & 시점 & 유형 & 직접 원인 & 사망자 \\
    \midrule
    챔플레인 타워 붕괴(미국) & 2021.6. & 붕괴 & 장기 구조 열화(철근 부식), 보수공사 지연 & 98명 \\
    튀르키예-시리아 지진 & 2023.2. & 지진 & M7.8 연쇄지진, 내진기준 미준수$\cdot$부실시공 & 약 59,000명+ \\
    마우이 라하이나 산불(미국) & 2023.8. & 산불 & 전력설비 발화 추정 + 침입종 초지$\cdot$강풍 & 102명 \\
    모로코 알하우즈 지진 & 2023.9. & 지진 & M6.8, 흙벽돌 조적조 밀집 산악마을 & 2,946명 \\
    리비아 데르나 홍수 & 2023.9. & 풍수해 & 폭풍 `다니엘' + 노후 댐 2기 연쇄 붕괴 & 4,352명+ \\
    노토반도 지진(일본) & 2024.1. & 지진 & M7.6, 노후 목조밀집$\cdot$반도부 접근성 제약 & 641명 \\
    LA 대형 산불(미국) & 2025.1. & 산불 & 샌타애나 강풍 + 장기 가뭄, 전력설비 발화 추정 & 30명 \\
    \bottomrule
  \end{tabularx}
  \par\vspace{8pt}
  {\scriptsize 재산피해: 튀르키예 직접피해 342억 달러+, 모로코 복구계획 120억 달러, 마우이 약 55억 달러, LA 경제손실 760억--1,310억 달러}
  \src{자료: 内閣府(2024), Fire Safety Research Institute(2024), NIST(2024), UCLA Anderson Forecast(2025), USGS(2023), World Bank(2023). 데르나는 실종 8천여 명 별도, 노토는 재해관련사 포함.}
\end{frame}

\begin{frame}{사례별 인명피해 규모}
  \centering
  \begin{tikzpicture}
    \begin{axis}[
      width=0.92\textwidth, height=0.80\textheight,
      xbar, bar shift=0pt, bar width=7pt,
      xmode=log, log basis x=10, xmin=3, xmax=400000,
      ytick={1,...,14}, yticklabels={광주 화정,부천 호텔,오송 지하차도,아리셀,LA 산불(미국),영남권 산불,챔플레인 타워(미국),마우이 산불(미국),이태원,무안 제주항공,노토반도(일본),모로코,데르나(리비아),튀르키예-시리아},
      y dir=reverse, ymin=0.4, ymax=14.6,
      yticklabel style={font=\scriptsize},
      xticklabel style={font=\scriptsize},
      xlabel={사망자 수 (명, 로그척도)}, xlabel style={font=\scriptsize},
      axis x line*=bottom, axis y line*=left, y axis line style={draw=none}, ytick style={draw=none},
      xmajorgrids, grid style={draw=black!12},
      point meta=explicit symbolic, nodes near coords,
      nodes near coords style={font=\tiny, text=black!75, anchor=west},
      area legend,
      legend style={at={(0.98,0.97)}, anchor=north east, font=\scriptsize, draw=none, fill=white},
      legend cell align=left,
    ]
      \addplot[fill=domestic, draw=none] coordinates {(6,1) [6] (7,2) [7] (14,3) [14] (23,4) [23] (32,6) [32] (159,9) [159] (179,10) [179]};
      \addplot[fill=foreign, draw=none] coordinates {(30,5) [30] (98,7) [98] (102,8) [102] (641,11) [641] (2946,12) [2,946] (4352,13) [4,352] (59000,14) [59,000]};
      \legend{국내, 해외}
    \end{axis}
  \end{tikzpicture}
  \src{저소득$\cdot$분쟁 지역의 피해가 압도적이나, 방재 선진국(미$\cdot$일$\cdot$한)에서도 관리 공백 영역에서는 세 자릿수 인명피해가 반복}
\end{frame}

\begin{frame}{유형별 분석 (1): 지진 $\cdot$ 화재 $\cdot$ 산불}
  \begin{block}{지진: 내진 규제의 집행 실패와 취약 주거의 방치}
    \small
    \begin{itemize}
      \item 튀르키예: `건축 사면(imar aff\fb{ı})'으로 기준 미달 건축물 수백만 동 합법화, 부실시공$\cdot$감리 부재
      \item 모로코: 내진 성능 없는 흙벽돌 조적조 / 노토: 과소$\cdot$고령지역 노후 목조, 재해관련사 장기화
      \item[$\Rightarrow$] 피해를 결정하는 것은 규제의 존재가 아니라 \alert{집행$\cdot$검증과 개보수의 실행}
    \end{itemize}
  \end{block}
  \begin{block}{화재$\cdot$산불: 기후변화와 관리 공백의 결합}
    \small
    \begin{itemize}
      \item 영남권: 일상적 실화 + 극한 기상, 고령자 대피체계 미작동, 침엽수 위주 산림구조
      \item 마우이$\cdot$LA: 전력설비 발화 + 방치된 연료(침입종 초지, WUI 확장) + 강풍$\cdot$가뭄
      \item[] \quad 보험사 인수 거절 확대라는 민간 리스크 신호가 \alert{예방투자로 연결되지 못함}
      \item 아리셀$\cdot$부천: 신종 위험물질, 스프링클러 소급 미적용 등 \alert{규제 사각지대}
    \end{itemize}
  \end{block}
\end{frame}

\begin{frame}{유형별 분석 (2): 붕괴 $\cdot$ 풍수해 $\cdot$ 사회재난}
  \begin{block}{붕괴: 품질$\cdot$유지관리 실패의 누적}
    \small
    \begin{itemize}
      \item 챔플레인 타워: 2018년 `중대한 구조적 손상' 보고 후 보수비 갈등으로 3년 이상 공사 지연
      \item 광주 화정: 무단 공법 변경, 동바리 조기 철거 $\to$ 2025 광명 신안산선 터널 붕괴로 반복
      \item[$\Rightarrow$] 위험이 감지 가능했으나 보수$\cdot$시공 품질을 강제$\cdot$지원할 \alert{제도적 장치 부재}
    \end{itemize}
  \end{block}
  \begin{block}{풍수해$\cdot$사회재난: 알려진 취약성과 경고의 방치}
    \small
    \begin{itemize}
      \item 데르나: 균열 위험을 경고한 논문까지 있었던 노후 댐을 분열된 행정이 방치
      \item 오송: 임시제방 부실 + 홍수경보 후 교통통제 미실시
      \item 이태원: 반복된 위험 신고가 예방적 통제로 이어지지 못함
      \item 무안: 부러지기 쉬워야(frangibility) 할 항행안전시설을 콘크리트로 시공
      \item[$\Rightarrow$] 위험 시나리오를 사전 검증하는 \alert{`레드팀형' 안전점검}의 부재
    \end{itemize}
  \end{block}
\end{frame}

\begin{frame}{유형별 원인구조 종합}
  \scriptsize
  \begin{tabularx}{\textwidth}{@{}l Y Y Y Y@{}}
    \toprule
    유형 & 직접 원인(트리거) & 구조적 원인 & 피해 특성 & 대표 사례 \\
    \midrule
    지진 & 단층 파열(자연력) & 내진기준 집행 실패, 취약주거 방치, 개보수 지연 & 대규모 인명피해, 재해관련사 장기화 & 튀르키예-시리아, 모로코, 노토 \\
    산불 & 실화$\cdot$전력설비 발화 + 강풍$\cdot$가뭄 & 연료 관리 부재, WUI 확장, 진화자원 한계 & 초광역 확산, 고령층 대피 실패 & 영남권, LA, 마우이 \\
    화재 & 위험물 열폭주, 전기적 발화 & 신종 위험물 규제 공백, 설비 소급 미비, 안전교육 부재 & 밀폐공간 다수 사망, 취약노동자 집중 & 아리셀, 부천 호텔 \\
    붕괴 & 구조 열화, 무단 공법 변경 & 유지보수 지연, 감리 형식화, 비용 갈등 & 순간 붕괴, 매몰 구조 장기화 & 챔플레인, 광주 화정 \\
    풍수해$\cdot$사회재난 & 극한강우, 군중밀집, 조류충돌 & 기반시설 노후 방치, 경고 무시, 시설 안전성 미검증 & 경고-대응 단절, 단시간 대량 인명피해 & 데르나, 오송, 이태원, 무안 \\
    \bottomrule
  \end{tabularx}
  \vspace{8pt}
  \begin{center}
    \small 유형을 불문한 공통 경로:\quad
    \fcolorbox{muted}{white}{위험의 사전 인지} $\to$
    \fcolorbox{accent}{white}{예방 투자$\cdot$집행 실패} $\to$
    \fcolorbox{muted}{white}{대응 단계 과부하}
  \end{center}
\end{frame}

\begin{frame}{사례조사의 함의}
  \begin{enumerate}
    \item \textbf{트리거 발생 $\neq$ 대형피해 `사건화'}
    \begin{itemize}\small
      \item 피해 규모는 자연적$\cdot$기술적 트리거보다, 그것이 대형피해로 번지는 것을 막는\\
            예방투자(내진보강, 유지보수, 연료관리, 소방설비, 대피체계)의 유무가 결정
      \item[$\Rightarrow$] 예방투자는 빈도뿐 아니라 \alert{심도와 손실분포의 꼬리(tail)}에 작용
    \end{itemize}
    \vspace{6pt}
    \item \textbf{볼록한(convex) 사회적 비용}
    \begin{itemize}\small
      \item 인명피해$\cdot$지역경제 충격$\cdot$사회적 갈등 등 직접 재산피해를 넘는 비용
      \item[$\Rightarrow$] 예방투자 평가에는 기대손실 감소뿐 아니라 \alert{대형손실 구간의 분포개선}이 반영되어야 함
    \end{itemize}
  \end{enumerate}
  \vspace{6pt}
  \begin{center}
    \small $\rightarrow$ 이론모형(빈도$\cdot$심도 동시 고려)과 실증분석(전체피해 vs.\ 우심피해 구분)의 출발점
  \end{center}
\end{frame}

% ==================================================================
\section{이론모형 개요}

\begin{frame}{이론모형의 핵심 (요약)}
  \begin{columns}[T]
    \begin{column}{0.5\textwidth}
      \textbf{복합 포아송 손실모형}
      \[
        S = \sum_{i=1}^{N(e_f)} q(e_s)\,X_i,\qquad N \sim \text{Poisson}\big(\lambda(e_f)\big)
      \]
      \small
      \begin{itemize}
        \item $e_f$: \textbf{빈도저감형} 투자 ($\lambda' < 0$)
        \item $e_s$: \textbf{심도저감형} 투자 ($q' < 0$)
        \item 사회적 목적함수에 볼록한 사회적 비용 $\Phi(S)$ 포함 ($\Phi''>0$)
        \item 실증용 함수형: 빈도 포아송/음이항,\\ 심도 감마/로그정규
      \end{itemize}
    \end{column}
    \begin{column}{0.48\textwidth}
      \textbf{실증분석에 필요한 함의}
      \small
      \begin{itemize}
        \item 같은 기대손실 감소라도 \alert{높은 손실구간}을 줄이는 투자의 사회적 가치가 더 큼
        \item 동일 기대손실 감소 하에서는 \alert{심도저감형}이 빈도저감형보다 사회비용 절감효과가 큼
        \item[$\Rightarrow$] 예방투자 효과는 \textbf{빈도/심도}, \textbf{전체/대형피해}로 나누어 보아야 함
      \end{itemize}
    \end{column}
  \end{columns}
  \src{※ 정리 및 증명은 본문 Ⅲ장 참조. 실증에서는 $e_f$, $e_s$를 구분한 예산자료가 없으므로, 총 예산계수를 축약형 탄력성으로 해석함.}
\end{frame}

% ==================================================================
\section{실증분석: 재난방재 예산과 자연재해 피해}

\begin{frame}{분석 개요}
  \begin{itemize}
    \item \textbf{분석 단위}: 17개 광역지방자치단체(시$\cdot$도) $\times$ 2015--2024년 (10개년)
    \item \textbf{설명변수}: 재난방재 예산 --- 배분 시차를 고려한 네 가지 사양
    \begin{center}\small
      \begin{tabular}{cccc}
        \toprule
        Current & Lag-1Y & Lag-2Y & Avg-3Y \\
        당해연도 & 1년 전 & 2년 전 & 직전 3개년 평균 \\
        \bottomrule
      \end{tabular}
    \end{center}
    \item \textbf{종속변수}: 손실분포의 구간별로 구분
    \begin{itemize}
      \item \textbf{전체} 자연재해 피해의 빈도와 심도
      \item \textbf{우심피해}(아주 심각한 대형피해)의 빈도와 심도
      \item 전체 피해 중 우심피해로 발전한 \textbf{비율}(우심피해 전환확률)
    \end{itemize}
    \item 전체/우심 구분 이유: 예방투자 효과가 \alert{손실분포의 어느 구간}에서 나타나는지가 사회적 가치 평가의 핵심이기 때문
  \end{itemize}
\end{frame}

\begin{frame}{자료의 구축}
  \scriptsize
  \begin{tabularx}{\textwidth}{@{}l l Y@{}}
    \toprule
    변수 & 정의 & 자료 원천 \\
    \midrule
    ExtCount / ExtLoss & 우심피해 건수 및 피해액 & 행정안전부 「재해연보 2024」 우심피해 내역 (전체 피해액의 약 70\%) \\
    AllCount & 전체 자연재해 피해 건수 & 「재해연보」(2015--2024) [기간별 피해발생 시도] \\
    AllLoss & 전체 자연재해 피해액 & 「재해연보 2024」 시도별$\cdot$연도별 피해현황 \\
    RecvBudget & 재난방재 예산 & 지방재정365: `재난방재/민방위' 분야 정책사업예산총계 \\
    Population, GRDP, Area & 인구, 지역내총생산, 면적 & KOSIS (DT\_1C91, 주민등록인구현황, 지적통계) \\
    PriorAllLoss / PriorExtLoss & 직전 기간 전체/우심 피해액 & 재해연보로부터 구축 (과거 피해 통제) \\
    \bottomrule
  \end{tabularx}
\end{frame}

\begin{frame}{추정 방법}
  \small
  \textbf{빈도 모형} --- 포아송 / 음이항(과산포 허용) 회귀
  \[
    \ln E[\text{Count}_{it}] = \gamma \ln(\text{Budget}_{i,t-k}) + \delta \ln(\text{PriorLoss}_{it}+1) + X_{it}'\beta + \mu_i + \tau_t
  \]
  \textbf{심도 모형} --- 건당 평균피해액(피해액 0 제외, 건수 가중)
  \begin{itemize}
    \item 로그링크 감마 회귀 \quad / \quad $\ln(\text{평균피해액})$에 대한 가우시안 회귀(로그정규 근사)
  \end{itemize}
  \textbf{우심피해 전환확률 모형} --- 이항 로짓(전체피해 건수 가중)
  \begin{itemize}
    \item 종속변수: ExtCount / AllCount $\Rightarrow$ 발생한 재해가 우심피해로 악화될 조건부 확률
  \end{itemize}
  \vspace{4pt}
  \begin{itemize}
    \item 통제변수: 과거 피해액, GRDP, 인구, 면적(로그) + \textbf{지역$\cdot$연도 고정효과}
    \item 표준오차: 시$\cdot$도 클러스터 강건 (df = 16)
    \item 예산계수 $\gamma$ = 해당 빈도/심도 지표에 대한 \alert{축약형 예산탄력성}
  \end{itemize}
\end{frame}

\begin{frame}{추정 결과: 재난방재 예산 계수}
  \scriptsize
  \centering
  \begin{tabular}{@{}l cccc@{}}
    \toprule
    종속변수 $\cdot$ 분포 & Current & Lag-1Y & Lag-2Y & Avg-3Y \\
    \midrule
    \rowcolor{rowhl}
    우심피해 빈도 $\cdot$ 포아송 & \textbf{$-$2.446} (0.032) & \textbf{$-$3.343} (0.005) & +0.768 (0.269) & $-$1.272 (0.289) \\
    \rowcolor{rowhl}
    우심피해 빈도 $\cdot$ 음이항 & $-$1.643 (0.090) & \textbf{$-$2.751} (0.027) & +0.164 (0.816) & $-$2.421 (0.073) \\
    우심피해 심도 $\cdot$ 감마 & $-$1.811 (0.066) & \textbf{+2.826} (0.001) & +1.486 (0.076) & +2.912 (0.098) \\
    우심피해 심도 $\cdot$ 로그노말 & $-$1.615 (0.102) & \textbf{+2.194} (0.012) & +1.220 (0.114) & +2.090 (0.179) \\
    \midrule
    전체피해 빈도 $\cdot$ 포아송 & +0.093 (0.616) & +0.263 (0.150) & +0.255 (0.056) & \textbf{+0.494} (0.038) \\
    전체피해 빈도 $\cdot$ 음이항 & +0.093 (0.616) & +0.263 (0.150) & +0.255 (0.056) & \textbf{+0.494} (0.038) \\
    전체피해 심도 $\cdot$ 감마 & \textbf{$-$1.495} (0.004) & $-$1.071 (0.142) & $-$0.871 (0.159) & \textbf{$-$2.513} (0.003) \\
    전체피해 심도 $\cdot$ 로그노말 & \textbf{$-$1.437} (0.017) & $-$0.512 (0.510) & $-$0.239 (0.731) & $-$1.273 (0.240) \\
    \midrule
    \rowcolor{rowhl}
    우심피해 전환확률 $\cdot$ 이항로짓 & \textbf{$-$3.350} (0.042) & \textbf{$-$4.215} (0.001) & +0.843 (0.401) & $-$2.748 (0.093) \\
    \bottomrule
  \end{tabular}
  \src{주: 괄호 안은 시도 클러스터 강건 표준오차(df=16)에 기초한 p-값, 굵은 글씨는 5\% 수준 유의. 심도 모형은 피해액 0 관측치 제외. 전체피해 빈도의 음이항 추정은 과산포 모수 추정의 반복 한계에 도달하여 포아송과 사실상 동일. 전환확률 모형은 과거 피해액 미포함(그 외 통제변수$\cdot$고정효과 동일).}
\end{frame}

\begin{frame}{결과의 해석 (1): 우심피해 --- 대형피해화의 억제}
  \begin{block}{우심피해 빈도: 분포 가정과 무관하게 음($-$)의 관계}
    \small
    \begin{itemize}
      \item Lag-1Y 기준: 포아송 $-$3.343 ($p$=0.005), 음이항 $-$2.751 ($p$=0.027)
      \item 방재예산이 큰 지역$\cdot$연도일수록 우심피해 건수가 유의하게 적음
      \item Lag-2Y에서만 부호가 반전되나 유의하지 않음
    \end{itemize}
  \end{block}
  \begin{block}{우심피해 전환확률: 당해$\cdot$전년도 예산이 악화 확률을 낮춤}
    \small
    \begin{itemize}
      \item Current $-$3.350 ($p$=0.042), Lag-1Y $-$4.215 ($p$=0.001), Avg-3Y $-$2.748 ($p$=0.093)
      \item 발생한 자연재해가 우심피해로 \alert{악화될 조건부 확률}이 낮음
      \item 시차별 부호 패턴이 우심피해 빈도 결과와 정합적
    \end{itemize}
  \end{block}
\end{frame}

\begin{frame}{결과의 해석 (2): 전체피해와 심도}
  \begin{itemize}
    \item \textbf{전체피해 빈도}: 모든 시차에서 양(+)이나 대체로 비유의 (Avg-3Y만 5\% 유의)
    \begin{itemize}\small
      \item[$\Rightarrow$] 방재예산이 자연재해 피해 사건의 \alert{발생 자체}를 줄인다는 증거는 없음
    \end{itemize}
    \item \textbf{전체피해 심도}: 모든 시차에서 음($-$) (감마$\cdot$로그노말 공통)
    \begin{itemize}\small
      \item 가장 뚜렷한 결과: 감마$\cdot$Avg-3Y $-$2.513 ($p$=0.003)
    \end{itemize}
    \item \textbf{우심피해 심도}: Current는 음(비유의), Lag-1Y는 유의한 양(+), 나머지는 대체로 비유의
    \item \textbf{통제변수}
    \begin{itemize}\small
      \item GRDP가 클수록 피해 빈도 증가 (노출효과)
      \item 과거 우심피해액이 클수록 이후 우심피해 빈도 감소 (사후 보강 학습효과)
      \item 연도효과는 2020년(전국적 집중호우) 전후에서 크게 나타남
    \end{itemize}
  \end{itemize}
\end{frame}

\begin{frame}{논의: 사례$\cdot$이론과의 대조}
  \centering
  \begin{tikzpicture}[
      box/.style={draw, rounded corners=3pt, align=center, minimum height=1.25cm, text width=3.6cm, font=\small},
      arr/.style={-{Stealth[length=6pt]}, thick, draw=black!60}]
    \node[box, draw=muted] (case) {\textbf{사례조사 (Ⅱ)}\\트리거와 대형피해\\`사건화'의 분리};
    \node[box, draw=muted, right=1.0cm of case] (theory) {\textbf{이론모형 (Ⅲ)}\\높은 손실구간 개선의\\사회적 가치가 더 큼};
    \node[box, draw=accent, right=1.0cm of theory] (emp) {\textbf{실증분석 (Ⅳ)}\\전체 빈도 무관,\\우심 빈도$\cdot$전환확률 감소};
    \draw[arr] (case) -- (theory);
    \draw[arr] (theory) -- (emp);
  \end{tikzpicture}
  \vspace{1.2em}

  \begin{alertblock}{종합 해석}
    방재예산은 자연현상의 \textbf{발생 자체}를 막기보다,\\
    발생한 재해가 \alert{일정 기준 이상의 대형피해로 발전하는 것을 억제}하는 경로로 작동할 수 있음
  \end{alertblock}
\end{frame}

\begin{frame}{식별의 한계}
  \begin{itemize}
    \item \textbf{내생성}: 방재예산은 과거 피해와 예상 위험에 반응하여 배분됨
    \begin{itemize}\small
      \item 과거 피해액 통제는 적절한 방향 / 추가로 지역별 강수량$\cdot$태풍 영향 등 당해연도 자연재해 충격 통제 검토
      \item 결과는 \alert{지역$\cdot$연도 고정효과와 관측 위험요인을 통제한 조건부 상관관계}로 해석
    \end{itemize}
    \vspace{4pt}
    \item \textbf{심도 결과의 해석 주의}
    \begin{itemize}\small
      \item Lag-1Y 예산과 우심피해 심도의 유의한 양(+)의 관계는 예산이 피해를 키운다는 의미가 아님
      \item 피해가 큰 지역에 다음 해 예산이 더 배분되는 \textbf{역인과}
      \item 우심피해 발생 관측치만을 조건으로 하는 \textbf{표본선택} 효과
    \end{itemize}
    \vspace{4pt}
    \item[$\Rightarrow$] \textbf{빈도 측 결과는 비교적 명확}하나, \textbf{심도효과는 안정적으로 식별되지 않음}
  \end{itemize}
\end{frame}

% ==================================================================
\section{결론 및 시사점}

\begin{frame}{결론}
  \begin{enumerate}
    \item \textbf{사례조사}: 대재해의 피해 규모는 트리거의 발생보다, 대형피해로의 사건화를 막는 \alert{예방투자의 존재 여부}와 밀접하게 관련
    \vspace{4pt}
    \item \textbf{이론모형}: 사회적 비용이 볼록할 때, 손실분포의 \alert{높은 구간}을 개선하는 예방투자(특히 심도저감형)의 사회적 가치가 더 큼
    \vspace{4pt}
    \item \textbf{실증분석}: 재난방재 예산은
    \begin{itemize}\small
      \item 전체 자연재해 피해의 발생 빈도와는 안정적인 음의 관계가 \textbf{없으나}
      \item 우심피해 빈도 및 우심피해로의 악화 확률과는 대체로 \textbf{음의 관계}
    \end{itemize}
  \end{enumerate}
  \vspace{6pt}
  \begin{center}
    $\Rightarrow$ 방재예산이 재해의 \alert{대형피해화를 억제}하는 경로로 작동한다는, 이론과 정합적인 잠정적 증거
  \end{center}
\end{frame}

\begin{frame}{시사점과 향후 과제}
  \begin{block}{정책적 시사점}
    \small
    \begin{itemize}
      \item 재난예방 성과를 전체 피해의 발생 빈도나 평균 피해규모만으로 평가해서는 안 됨
      \item 우심피해처럼 \alert{손실분포의 높은 구간}에서 나타나는 피해를 구분하여 평가할 필요
      \item 규제의 존재보다 \textbf{집행$\cdot$검증$\cdot$유지보수}에 대한 예방투자가 대형화를 좌우
    \end{itemize}
  \end{block}
  \begin{block}{연구의 한계와 향후 과제}
    \small
    \begin{itemize}
      \item 제한된 지역 패널에서 확인된 \textbf{조건부 관계} --- 인과효과나 개별 재난의 대형화 과정을 직접 식별한 것은 아님
      \item 빈도저감형/심도저감형 예산을 구분한 세분화된 예방사업 자료와 재난충격 자료가 축적되면 경로를 직접 분석 가능
    \end{itemize}
  \end{block}
\end{frame}

\begin{frame}{주요 참고문헌}
  \scriptsize
  \begin{itemize}
    \item Ehrlich, I., \& Becker, G. S. (1972). Market Insurance, Self-Insurance, and Self-Protection. \textit{Journal of Political Economy}, 80(4), 623--648.
    \item Lee, K. (1998). Risk Aversion and Self-Insurance-cum-Protection. \textit{Journal of Risk and Uncertainty}, 17(2), 139--150.
    \item Wong, K. P. (2016). Precautionary Self-Insurance-cum-Protection. \textit{Economics Letters}, 145, 152--156.
    \item 홍지민 (2022). 손실통제: 2기간 모형에서 비교정학분석을 중심으로. 보험학회지 131, 1--22.
    \item 배은경, 김소연, 김창기 (2018). 극단치 이론을 이용한 지진위험 모델링과 관리. 리스크관리연구 29(2), 1--24.
    \item 이동근, 김호걸, 남상욱 (2017). 기후변화 리스크의 위험성과 확장성 분석. 리스크관리연구 28(2), 69--94.
    \item 행정안전부. 「재해연보 2024」 (2025).
    \item World Bank (2023). GRADE Report: Kahramanmara\fb{ş} Earthquakes -- T\fb{ü}rkiye.
    \item NIST (2024). Champlain Towers South Partial Collapse: Investigation Update.
    \item 지방재정365 재정 데이터셋 / 통계청 KOSIS
  \end{itemize}
\end{frame}

\begin{frame}{맺음말: 성사재천 모사재인}
  \centering
  {\small\color{muted} 고사: 謀事在人 成事在天 (모사재인 성사재천) --- 일을 꾸미는 것은 사람이나, 이루는 것은 하늘에 달려 있다}

  \vspace{1.4em}
  {\Huge\bfseries 成事在天\quad\alert{謀事在人}}\par
  \vspace{0.4em}
  {\large 성사재천 \quad 모사재인}

  \vspace{1.6em}
  \begin{columns}[T]
    \begin{column}{0.44\textwidth}
      \begin{block}{成事在天}
        \small 자연재해의 \textbf{발생 자체}는 인간의 힘으로 막을 수 없다\\[3pt]
        {\footnotesize\color{muted} 전체 자연재해 피해 빈도는 방재예산과 뚜렷한 관계 없음}
      \end{block}
    \end{column}
    \begin{column}{0.44\textwidth}
      \begin{alertblock}{謀事在人}
        \small 그 피해가 \alert{지나치게 커지지 않도록 통제}하는 것이 인간의 역할이다\\[3pt]
        {\footnotesize\color{muted} 우심피해 빈도와 우심피해 전환확률은 방재예산과 음($-$)의 관계}
      \end{alertblock}
    \end{column}
  \end{columns}
\end{frame}

\begin{frame}[standout]
  감사합니다
\end{frame}

\end{document}
